\begin{frame}{자주 하는 실수: ``최신 모델''부터 잡기}
  \begin{itemize}
    \item 가장 흔한 실패 패턴: Transformer(Week 12)·LLM(Week 17)이
      ``산업에서 쓰는 것''이라는 이유로 \textbf{Block A를 건너뛰고
      Week 12로 직행}.
    \item 실제 경험: 작은 네트워크의 손실이 안 떨어지는데
      (a) 데이터 스케일? (b) 학습률? (c) 구조 자체의 문제인지 구분 못
      하고 ``일단 층을 늘려보자''로 끝남.
    \item 한 달 뒤엔 \textbf{원인 모를 손실 곡선}만 남는다.
  \end{itemize}
  \vskip0.5em
  \begin{block}{실전에서 막히는 진짜 이유}
    Week 6의 ``train/val/test 분리''와 Week 2의 ``손실함수를 손으로
    한 번 계산'' 같은 \textbf{기초 도구가 없기} 때문이다.
  \end{block}
\end{frame}
